\documentclass[a4paper]{article}
\usepackage{amsmath}
\usepackage{amsfonts}
\usepackage[affil-it]{authblk}
\usepackage[backend=bibtex,style=numeric]{biblatex}

\usepackage{geometry}
\geometry{margin=1.5cm, vmargin={0pt,1cm}}
\setlength{\topmargin}{-1cm}
\setlength{\paperheight}{29.7cm}
\setlength{\textheight}{25.3cm}
\setlength{\parindent}{0pt}

\begin{document}
% =================================================
\title{Report For Solving Nonlinear Equations}

\author{WangHao 3220103613
  \thanks{Electronic address: \texttt{3220103613@zju.edu.cn}}}
\affil{(Information and Computing Sciences, Class 2201), Zhejiang University }


\date{Due time: \today}

\maketitle






% ============================================
\section*{A.}
My header file is written in the directory \texttt{.\textbackslash lib}.

In \texttt{Function.hpp}, I wrote a class called \texttt{Function}. 
This is an abstract class that defines the characteristics of the function class: 
its operations include evaluating the function value and the derivative value. 
These operations can be override. 
This abstract class can be inherited to generate different function classes.

In \texttt{EquationSolver.hpp}, I wrote a class called \texttt{EquationSolver}. 
This is an abstract class that defines the characteristics of an equation solver for \( F(x) = 0 \): 
its elements are functions (references) \( F \), and its operations include finding roots. 
This operation can be override. 
This abstract class can be inherited to generate different function solvers. 
In this header file, this abstract class is inherited to create three function solvers, 
representing three different methods.

\section*{B.}
Run B, we have:


Solving $x^{-1} - \tan x$ on $[0, \pi/2]$

A root is: 0.860333

Solving $x^{-1}-2^{x}$ on $[0,1]$

A root is: 0.641186

Solving $2^{-x}+e^{x}+2cosx-6$ on $[1,3]$

A root is: 1.82938

Solving $\frac{x^{3}+4x^{2}+3x+5}{2x^{3}-9x^{2}+18x-2}$ on $[0,4]$

The function is not continuous, the root is probably not correct. 
A root is: 0.117877


In function4, F(root) is huge, thus the program throw the hint:
The function is not continuous, the root is probably not correct. 

\section*{C.}
Run C, we have:


Solving $x-\tan(x)$ near 4.5

A root is 4.49341

Solving $x-\tan(x)$ near 7.7

A root is 7.72525


We have found root close to the initial value successfully.

\section*{D.}
Run D, we have:


Solving $\sin(x/2)-1$ with $x_0=0, x_1=1.5708$

A root is 3.14093

Solving $e^{x}-\tan(x)$ with $x_0=1, x_1=1.4$

A root is 1.30633

Solving $x^{3}-12x^{2}+3x+1$ with $x_0=0, x_1=-0.5$

A root is -0.188685


Take another set of initial values.

Solving $\sin(x/2)-1$ with $x_0=0.261799, x_1=0.241661$
A root is 3.14074

We could find that two roots of F1 are different. There are 2 close roots.

Solving $e^{x}-\tan(x)$ with $x_0=1.1, x_1=1.3$

A root is 1.30633
Solving $x^{3}-12x^{2}+3x+1$ with $x_0=-0.1, x_1=-0.3$

A root is -0.188685

The roots of F2 or F3 are same with different initial value (I've tried a lot.). 
Thus F2 or F3 may not have roots close to roots solved above.

\section*{E.}
Run E, we have:


Solving Function with $a=0, b=1$

A root is 0.164062

Solving Function with $x_0=0$

A root is 0.166166

Solving Function with $x_0=0, x_1=1$

A root is 0.166166


And we can find that the root is more close to 0.166166 than 0.164062. 
With the same precision 0.01ft, Bisection Method converges slower than Newton Method and Secant Method.
This is in line with the theory.

\section*{F.}
Run F, we have:

\subsection*{(a)}

F.(a)
Solving Function with $x_0=33$

A root is 32.9722

We have found root close to the initial value 33 successfully.

\subsection*{(b)}
F.(b)

Solving Function with $x_0=33$

A root is 33.1689

The root varies with the change in parameter $D$.

\subsection*{(c)}
F.(c)

Solving Function with $x_0=33, x_1=100000$

A root is 392.972

Solving Function with $x_0=33, x_1=200000$

A root is 447873

Solving Function with $x_0=33, x_1=300000$

A root is 10947


As two initial value are so far away, secant can't fitting tangent line now. 
After first iteration, the iteration point went so far away the initial point. 
So we have solved the root far away 33.


\end{document}